\section{Proofs of nonparametric convergence results}
\label{sec:proof-nonparametric}

\newcommand{\noise}{\xi}
\newcommand{\localcomplexity}{\mc{R}}


In this technical appendix, we include proofs of the results from
Section~\ref{sec:semiparametric} as well as a few additional results, which
are essentially corollaries of results on localized complexities and
nonparametric regression models, though we require a few modifications
because our setting is slightly non-standard.


\subsection{Preliminary results}

To set notation and to make reading it self-contained, we provide some
definitions. The $L^r(P)$ norm of a function or random vector is
$\norm{f}_{L^r(P)} = (\int |f|^r dP)^{1/r}$, so that its $L^2(P_n)$-norm is
$\LtwoPn{f}^2 = \frac{1}{n} \sum_{i = 1}^n f(X_i)^2$.  We consider the
following abstract nonparametric regression setting: we have a function
class $\mc{F} \subset \{\R \to \R\}$ with $f\opt \in \mc{F}$, and our
observations follow the model
\begin{equation}
  \label{eqn:nonparametric-model}
  Y_i = f\opt(X_i) + \noise_i,
\end{equation}
but instead of observing $(X_i, Y_i)$ pairs we observe $(\wt{X}_i, Y_i)$
pairs, where $\wt{X}_i$ may not be identical to $X_i$ (these play the roll
of $\<u\opt, X_i\>$ versus $\<\uinit, X_i\>$ in the results to come).  We
assume that $\noise_i$ are bounded so that $\sup \noise - \inf \noise \le
1$, independent, and satisfy the conditional mean-zero property that
$\E[\noise_i \mid X_i] = 0$ (though $\noise_i$ may not be independent of
$X_i$).  For a (thus far unspecified) function class $\mc{F}$ we set
\begin{equation*}
  \what{f} = \argmin_{f \in \mc{F}} P_n(Y - f(\wt{X}))^2.
\end{equation*}

We now demonstrate that the error $\LtwoPns{\what{f} - f\opt}^2 =
\frac{1}{n} \sum_{i = 1}^n (\what{f}(X_i) - f\opt(X_i))^2$ can be bounded by
a combination of the errors $\wt{X}_i - X_i$ and local complexities of the
function class $\mc{F}$.  Our starting point is a local complexity bound
analagous to the localization results available for in-sample prediction
error in nonparametric regression~\cite[cf.][Thm.~13.5]{Wainwright19}.  To
present the results, for a function class $\mc{H}$ we define the
localized $\noise$-complexity
\begin{equation*}
  \localcomplexity_n(u; \mc{H})
  \defeq \E\left[ \sup_{h \in \mc{H},
      \LtwoPns{h} \le u}
    \left|n^{-1} \sum_{i = 1}^n \noise_i h(x_i) \right| \right],
\end{equation*}
where we treat the expectation conditionally on $X_i$ and $\noise_i$ are
random (i.e.\ $\noise_i = Y_i - f\opt(X_i)$).  For the
model~\eqref{eqn:nonparametric-model}, we define the centered class
$\mc{F}\opt = \{f - f\opt \mid f \in \mc{F}\}$, which is star-shaped as
$\mc{F}$ is a convex set.\footnote{ A set $\mc{H}$ is \emph{star-shaped} if
  for all $h \in \mc{H}$, if $\alpha \in [0, 1]$ then $\alpha h \in
  \mc{H}$.}  We say that $\delta$ satisfies the \emph{critical radius
  inequality} if
\begin{equation}
  \frac{1}{\delta} \localcomplexity_n(\delta; \mc{F}\opt)
  \le \delta.
  \label{eqn:critical-inequality}
\end{equation}
With this, we can provide a proposition giving a high-probability
bound on the in-sample prediction
error of the empirical estimator $\what{f}$, which is essentially
identical to~\cite[Thm.~13.5]{Wainwright19}, though we require a few
modifications to address that we observe $\wt{X}_i$ and not $X_i$
and that the noise $\noise_i$ are bounded but not Gaussian.
\begin{proposition}
  \label{proposition:in-sample-convergence}
  Let $\mc{F}$ be a convex function class, $\delta_n > 0$ satisfy the
  critical inequality~\eqref{eqn:critical-inequality},
  and let
  $\gamma^2 = \sup_{f \in \mc{F}}
  \frac{1}{n} \sum_{i = 1}^n (f(X_i) - f(\wt{X}_i))^2$.
  Then for $t \ge \delta_n$,
  \begin{equation*}
    \P\left(\LtwoPns{\what{f} - f\opt}^2 \ge 30 t \delta_n
    + 25 \gamma \max\{\gamma, \LtwoPns{\noise}\}
    \mid X_1^n, \wt{X}_1^n
    \right)
    \le \exp\left(-\frac{n t \delta_n}{4}\right).
  \end{equation*}
\end{proposition}
\noindent
See Section~\ref{sec:proof-in-sample-convergence} for a proof of the
proposition.

Revisiting the critical radius inequality~\eqref{eqn:critical-inequality},
we can also apply \cite[Corollary~13.7]{Wainwright19}, which allows
us to use an entropy integral to guarantee bounds on the critical
radius. Here, we again fix $X_1^n = x_1^n$, and for a function
class $\mc{H}$ we let
$\mc{B}_n(\delta; \mc{H}) = \{h \in \starhull(\mc{H})
\mid \LtwoPn{h} \le \delta\}$. Let $N_n(t; \mc{B})$ denote the
$t$-covering number of $\mc{B}$ in $\LtwoPn{n\cdot}$-norm. Then
modifying a few numerical constants, we have
\begin{corollary}[Wainwright~\cite{Wainwright19}, Corollary 13.7]
  \label{corollary:localization-via-dudley}
  Let the conditions of
  Proposition~\ref{proposition:in-sample-convergence} hold. Then
  for a numerical constant $C \le 16$,
  any $\delta \in [0, 1]$ satisfying
  \begin{equation*}
    \frac{C}{\sqrt{n}} \int_{\delta^2 / 4}^\delta
    \sqrt{\log N_n(t; \mc{B}_n(\delta, \mc{F}\opt))} dt
    \le \frac{\delta^2}{2}
  \end{equation*}
  satisfies the critical inequality~\eqref{eqn:critical-inequality}.
\end{corollary}
\noindent
As an immediate consequence of this inequality, we have the following:
\begin{corollary}
  \label{corollary:lipschitz-criticality}
  Assume $|x_i| \le b$ for all $i \in [n]$ and that
  $\mc{F}$ is contained in the class of $M$-Lipschitz functions
  with $f(0) = 0$ (or any other fixed constant).
  Then for a numerical
  constant $c < \infty$, the choice
  \begin{equation*}
    \delta_n = c \left(\frac{M b}{n}\right)^{1/3}
  \end{equation*}
  satisfies the critical inequality~\eqref{eqn:critical-inequality}.
\end{corollary}
\begin{proof}
  The covering numbers $N_\infty$
  for the class $\mc{F}$ of $M$-Lipschitz functions
  on $[-b, b]$
  satisfying $f(0) = 0$
  in supremum norm $\linf{f} = \sup_{x \in [-b, b]} |f(x)|$ satisfy
  $\log N_\infty(t; \mc{F}) \lesssim \frac{M b}{t}$
  \cite[cf.][Example 5.10]{Wainwright19}. Using that $N_n \le N_\infty$,
  we thus have
  \begin{equation*}
    \int_{\delta^2/4}^\delta
    \sqrt{\log N_n(t; \mc{B}_n(\delta, \mc{F}\opt))}
    \lesssim
    \int_{\delta^2/4}^\delta
    \sqrt{\frac{Mb}{t}} dt
    = 2 \sqrt{Mb} \left(\sqrt{\delta} - \delta/4\right)
    \le 2 \sqrt{Mb \delta}
  \end{equation*}
  whenever $\delta \le 1$.
  For a numerical constant $c > 0$ it suffices in
  Corollary~\ref{corollary:localization-via-dudley} to choose
  $\delta$ satisfying
  $c \frac{1}{\sqrt{n}} \sqrt{Mb \delta} \le \delta^2$, or
  $\delta = c (\frac{Mb}{n})^{1/3}$.
\end{proof}


\subsection{Proof of Proposition~\ref{proposition:consistency-of-sigma}}
\label{sec:proof-consistency-of-sigma}

We assume without loss of generality that $m = 1$, as nothing in the
proof changes except notationally (as we assume $m$ is fixed).
We apply Proposition~\ref{proposition:in-sample-convergence} and
Corollary~\ref{corollary:lipschitz-criticality}.  For
notational simplicity, let $Q$ denote the measure
on $\R$ that $Y \<u\opt, X\>$ induces for $X \sim P$, and $Q_n$ similarly for
$P_n$.
We first show that
$\LtwoQns{\sigma_n - \sigma\subopt}$ converges.  First, we
recall~\cite[Lemma 3]{Owen90} that $\max_{i \le n} n^{-1/k} \ltwo{X_i} \cas
0$ as $n \to \infty$. Thus there is a (random) $B < \infty$ such that
$\max_{i \le n} \ltwo{X_i} \le B n^{1/k}$ for all $n$. Therefore,
Corollary~\ref{corollary:lipschitz-criticality} implies that the choice
$\delta_n = c (\frac{\lipconst B}{n^{1-1/k}})^{1/3}$ satisfies the critical
inequality~\eqref{eqn:critical-inequality}, and taking $\gamma^2_n =
\frac{\lipconst^2}{n} \sum_{i = 1}^n \<\uinit - u\opt, X_i\>^2 \le M^2 \uerror_n^2
\LtwoPn{X}^2$, we have that for $t \ge \delta_n$,
\begin{equation*}
  \LtwoPns{\sigma_n - \sigma\subopt}^2
  \lesssim t \delta_n + \lipconst^2 \uerror_n^2 \LtwoPn{X}^2
  ~~~ \mbox{with~probability~at~least}~
  1 - \exp\left(-\frac{n t \delta_n}{4}\right)
\end{equation*}
on the event that $\max_{i \le n} \ltwo{X_i} \le B n^{1/k}$ for all $n$,
where we have conditioned (in the probability) on $X_i$.
As $\LtwoPn{X}^2 \cas \E[\ltwo{X}^2]$, we may choose
$t = \delta_n \gg 1/n^{1/3}$ and find that the Borell-Cantelli lemma
then implies that with probability $1$, there is a random $C < \infty$
such that
\begin{equation}
  \label{eqn:ltwopn-sigma-converge}
  \LtwoQn{\sigma_n - \sigma\opt}^2 \le C \left(
  n^{\frac{2}{3k} - \frac{2}{3}} + \epsilon_n^2 \right)
\end{equation}
for all $n$ with probability $1$.

Finally we argue that $\LtwoQ{\sigma_n - \sigma\opt} \cas 0$.
Let $b < \infty$ be otherwise arbitrary, and let
$\mc{G}_b$ be the collection of $2\lipconst$-Lipschitz functions on
$[-b, b]$ with $\linf{g} \le 1$ and $g(0)$ for $g \in \mc{G}_b$, noting
that $\sigma_n - \sigma\opt$ restricted to $[-b, b]$ evidently
belongs to $\mc{G}_b$. Then we have
\begin{align}
  \lefteqn{\LtwoQ{\sigma_n - \sigma\subopt}^2} \nonumber \\
  & = \LtwoQn{\sigma_n - \sigma\subopt}^2
  + \int_{|t| \le b} (\sigma_n(t) - \sigma\subopt(t))^2
  (dQ(t) - dQ_n(t))
  + \int_{|t| > b} (\sigma_n(t) - \sigma\subopt(t))^2 
  (dQ(t) - dQ_n(t)) \nonumber \\
  & \le \LtwoQn{\sigma_n - \sigma\subopt}^2
  + \sup_{g \in \mc{G}_b} |Q g^2 - Q_n g^2|
  + Q([-b, b]^c) + Q_n([-b, b]^c),
  \label{eqn:three-terms-bike}
\end{align}
where the inequality used that $\linf{\sigma_n - \sigma\subopt} \le 1$
by construction.
The first term in inequality~\eqref{eqn:three-terms-bike} we have
already controlled. We may control the second supremum term
almost immediately using Dudley's entropy integral
and a Rademacher contraction inequality.
Indeed, we have
\begin{equation*}
  \E[\sup_{g \in \mc{G}_b} |Q_n g^2 - Q g^2|]
  \stackrel{(i)}{\le} 2 \E[\sup_{g \in \mc{G}_b} |Q_n^0 g^2|]
  \stackrel{(ii)}{\le} 4 \E[\sup_{g \in \mc{G}_b} |Q_n^0|]
  \stackrel{(iii)}{\lesssim} \frac{1}{\sqrt{n}} \int_0^\infty
  \sqrt{\log N_\infty(t; \mc{G}_b)} dt,
\end{equation*}
where inequality~$(i)$ is a standard symmetrization inequality, $(ii)$ is
the Rademacher contraction inequality~\cite[Ch.~4]{LedouxTa91} applied to
the function $t \mapsto t^2$, which is $2$ Lipschitz for $t \in [-1, 1]$,
and $(iii)$ is Dudley's entropy integral bound. As the $\sup$-norm
log-covering numbers of $\lipconst$-Lipschitz functions on $[-b, b]$ scale
as $\frac{\lipconst b}{t}$ for $t \le \lipconst b$ and are 0 otherwise, we
obtain $\E[\sup_{g \in \mc{G}_b} |Q_n g^2 - Q g^2|] \lesssim
\frac{\lipconst b}{\sqrt{n}}$.  The bounded-differences inequality then implies that
for any $t > 0$,
\begin{equation*}
  \P\left(\sup_{g \in \mc{G}_b} |(Q_n - Q)g^2|
  \ge c \frac{\lipconst b}{\sqrt{n}} + t \right) \le
  \P\left(\sup_{g \in \mc{G}_b} |(Q_n - Q) g^2|
  \ge \E[\sup_{g \in \mc{G}_b} |(Q_n - Q) g^2|] + t \right)
  \le \exp(-c n t^2).
\end{equation*}
Finally, the final term in the bound~\eqref{eqn:three-terms-bike}
evidently satisfies $Q([-b, b]^c) \le \frac{\E[\ltwo{X}^k]}{b^k}$
and $\sup_b |Q_n([-b,b]^c) - Q([-b, b])| \le 2 \sqrt{t / n}$
with probability at least $1 - e^{-2t^2}$ by the DKW inequality.
We thus find by the Borel-Cantelli lemma that simultaneously
for all $b < \infty$, with probability at least
$1 - e^{-nt^2}$ we have
\begin{equation*}
  \LtwoQ{\sigma_n - \sigma\subopt}^2 \le
  \LtwoQn{\sigma_n - \sigma\subopt}^2
  + \frac{C \lipconst b}{\sqrt{n}} + C t + \frac{\E[\ltwo{X}^k]}{b^k},
\end{equation*}
where $C$ is a numerical constant.

Substituting inequality~\eqref{eqn:ltwopn-sigma-converge} into
the preceding display and taking $b = n^{-\frac{1}{2(k + 1)}}$, we
get the result.

\subsection{Proof of Proposition~\ref{proposition:in-sample-convergence}}
\label{sec:proof-in-sample-convergence}

We begin with an extension of the
familiar basic inequality~\cite[e.g.][Eq.~(13.18)]{Wainwright19}, where
we see by convexity that for any $\eta > 0$ we have
\begin{align*}
  \sum_{i = 1}^n (Y_i - \what{f}(X_i))^2
  & \le (1 + \eta) \sum_{i = 1}^n (Y_i - \what{f}(\wt{X}_i))^2
  + (1 + 1/\eta) \sum_{i = 1}^n (\what{f}(\wt{X}_i) - \what{f}(X_i))^2 \\
  & \le (1 + \eta) \sum_{i = 1}^n (Y_i - f\opt(\wt{X}_i))^2
  + (1 + 1/\eta) \sum_{i = 1}^n (\what{f}(\wt{X}_i) - \what{f}(X_i))^2 \\
  & \le (1 + \eta)^2 \sum_{i = 1}^n (Y_i - f\opt(X_i))^2
  + (2 + \eta + 1/\eta) \sum_{i = 1}^n
  \left[(\what{f}(\wt{X}_i) - \what{f}(X_i))^2
    + (f\opt(\wt{X}_i) - f\opt(X_i))^2 \right].
\end{align*}
Noting that $Y_i = f\opt(X_i) + \noise_i$, algebraic manipulations yield
that if
$\Delta = [f\opt(X_i) - \what{f}(X_i)]_{i = 1}^n$ is the error vector, then
\begin{equation*}
  \LtwoPn{\Delta}^2 - \frac{2}{n} \noise^T \Delta
  + \LtwoPn{\noise}^2 \le
  (1 + \eta)^2 \LtwoPn{\noise}^2
  + \frac{2 + \eta + 1/\eta}{n} \sum_{i = 1}^n
  \left[(\what{f}(\wt{X}_i) - \what{f}(X_i))^2
    + (f\opt(\wt{X}_i) - f\opt(X_i))^2 \right].
\end{equation*}
Simplifying, we obtain the following:
\begin{lemma}
  \label{lemma:basic-inequality}
  Let $\gamma^2 = \sup_{f \in \mc{F}}
  \frac{1}{n} \sum_{i = 1}^n (f(X_i) - f(\wt{X}_i))^2$
  and $\Delta = [\what{f}(X_i) - f\opt(X_i)]_{i = 1}^n$.
  Then
  \begin{align*}
    \LtwoPn{\Delta}^2
    & \le \inf_\eta \left\{(2 \eta + \eta^2) \LtwoPn{\noise}^2
    + \frac{2}{n} \noise^T \Delta
    + (4 + 2 \eta + 2/\eta) \gamma^2 \right\} \\
    & \le \frac{2}{n} \noise^T \Delta
    + 11 \gamma \max\left\{\gamma, \LtwoPn{\noise}\right\}.
  \end{align*}
\end{lemma}
\begin{proof}
  The first inequality is algebraic manipulations and
  uses that our choice of $\eta$ was arbitrary.
  For the second, we consider two cases:
  that $\LtwoPn{\noise} \ge \gamma$ and that $\LtwoPn{\noise} \le \gamma$.
  In the first case, we consider $\eta \le 1$,
  yielding the simplified bound
  that
  $\LtwoPn{\Delta}^2 \le \frac{2}{n} \noise^T \Delta
  + 3 \eta \LtwoPn{\noise}^2 + \frac{8 \gamma^2}{\eta}$ for
  $\eta \le 1$. Taking $\eta = \gamma / \LtwoPn{\noise}$ gives
  that
  \begin{equation*}
    \LtwoPn{\Delta}^2
    \le \frac{2}{n} \noise^T \Delta
    + 11 \LtwoPn{\noise} \gamma.
  \end{equation*}
  In the case that $\gamma \ge \LtwoPn{\noise}$, we choose $\eta = 1$,
  and the bound simplifies to
  $\LtwoPn{\Delta}^2 \le \frac{2}{n} \noise^T \Delta
  + 3 \LtwoPn{\noise}^2 + 8 \gamma^2
  \le \frac{2}{n} \noise^T \Delta + 11 \gamma^2$.
\end{proof}

We now return to the proof of the proposition proper.
We begin with an essentially immediate extension of the result~\cite[Lemma
  13.12]{Wainwright19}.  We let $\mc{H}$ be an arbitrary star-shaped
function class.
Define the event
\begin{equation}
  \label{eqn:A-wainwright-event}
  A(u) \defeq \left\{
  \mbox{there~exists~} g \in \mc{H} ~ \mbox{s.t.}~
  \LtwoPn{g} \ge u ~ \mbox{and} ~
  \left|P_n \noise g\right| \ge 2 \LtwoPn{g} u
  \right\},
\end{equation}
which we treat conditionally on $X_1^n = x_1^n$ as in the definition
of the local complexity $\localcomplexity_n$.
(Here the noise $\noise_i$ are still random, taken conditionally
on $X_1^n = x_1^n$.)

\begin{lemma}[Modification of Lemma 13.12 of~\citet{Wainwright19}]
  \label{lemma:star-shaped-localization}
  Let $\mc{H}$ be a start-shaped function class
  and let $\delta_n > 0$ satisfy the critical radius inequality
  \begin{equation*}
    \frac{1}{\delta} \localcomplexity_n(\delta; \mc{H}) \le
    \delta.
  \end{equation*}
  Then for all $u \ge \delta_n$, we have
  \begin{equation*}
    \P(A(u)) \le \exp\left(-\frac{n u^2}{4} \right).
  \end{equation*}
\end{lemma}
\noindent
Deferring the proof of Lemma~\ref{lemma:star-shaped-localization}
to Section~\ref{sec:proof-star-shaped-localization},
we can parallel the argument for~\cite[Thm.~13.5]{Wainwright19}
to obtain our proposition.

Let $\mc{H} = \mc{F}\opt$ in Lemma~\ref{lemma:star-shaped-localization}.
Whenever $t \ge \delta_n$, we have
$\P(A(\sqrt{t \delta_n})) \le e^{-nt \delta_n / 4}$. We consider the two
cases that $\LtwoPns{\Delta} \lessgtr \sqrt{t \delta_n}$. In the
former that $\LtwoPns{\Delta} \le \sqrt{t \delta_n}$, we have
nothing to do. In the latter, we have $\what{f} - f\opt \in \mc{F}\opt$
while $\LtwoPns{\Delta} > \sqrt{t \delta_n}$, so that if
$A(\sqrt{t \delta_n})$ fails then we must have
\begin{equation*}
  \left|\frac{1}{n} \sum_{i = 1}^n \noise_i \Delta(x_i)\right|
  \le 2 \LtwoPn{\Delta} \sqrt{t \delta_n}.
\end{equation*}
From the extension of the basic inequality in
Lemma~\ref{lemma:basic-inequality} we see that
\begin{equation*}
  \LtwoPn{\Delta}^2 \le 
  4 \LtwoPn{\Delta} \sqrt{t \delta_n} +
  11 \max\left\{\gamma^2, \gamma \LtwoPn{\noise}\right\}.
\end{equation*}
Solving for $\LtwoPn{\Delta}$ then yields
\begin{equation*}
  \LtwoPn{\Delta}
  \le \frac{4 \sqrt{t \delta_n} + \sqrt{16 t \delta_n
      + 44 \gamma \max\{\gamma, \LtwoPn{\noise}\}}}{2}
  \le 4 \sqrt{t \delta_n} +
  \sqrt{11 \gamma \max\{\gamma, \LtwoPn{\noise}\}}.
\end{equation*}
Simplifying gives Proposition~\ref{proposition:in-sample-convergence}.

\subsubsection{Proof of Lemma~\ref{lemma:star-shaped-localization}}
\label{sec:proof-star-shaped-localization}

Mimicking the proof of~\cite[Lemma 13.12]{Wainwright19}, we
begin with~\cite[Eq.~(13.40)]{Wainwright19}:
\begin{equation*}
  \P(A(u))
  \le \P(Z_n(u) \ge 2 u^2)
  ~~ \mbox{for} ~~
  Z_n(u) \defeq \sup_{g \in \mc{H},
    \LtwoPns{g} \le u} \left|n^{-1} \sum_{i=1}^n \noise_i g(x_i)\right|.
\end{equation*}
Now note that if $\LtwoPns{g} \le u$, then the function $\noise \mapsto
|n^{-1} \sum_{i = 1}^n \noise_i g(x_i)|$ is $u / \sqrt{n}$-Lipschitz
with respect to the $\ell_2$-norm, so that convex concentration
inequalities~\cite[e.g.][Theorem 3.4]{Wainwright19}
imply that
$\P(Z_n(u) \ge \E[Z_n(u)] + t)
\le \exp(-\frac{t^2 n}{4 b^2 u^2})$
whenever $\sup \noise - \inf \noise \le b$, and so
for $b = 1$ we have
\begin{equation*}
  \P(Z_n(u) \ge \E[Z_n(u)] + u^2) \le
  \exp\left(-\frac{n u^2}{4}\right).
\end{equation*}
As $\E[Z_n(u)] = \localcomplexity_n(u)$, we
finally use that the normalized complexity $t \mapsto
\frac{\localcomplexity_n(t)}{t}$ is non-decreasing~\cite[Lemma
  13.6]{Wainwright19} to obtain that
for $u \ge \delta_n$,
$\frac{1}{u} \localcomplexity_n(u)
\le \frac{1}{\delta_n} \localcomplexity_n(\delta_n)
\le \delta_n$, the last inequality by assumption.
In particular, we find that for $u \ge \delta_n$ we have
$\E[Z_n(u)] = \localcomplexity_n(u)
\le u \delta_n \le u^2$, and so
\begin{equation*}
  \P(Z_n(u) \ge 2 u^2) \le
  \P(Z_n(u) \ge \E[Z_n(u)] + u^2)
  \le \exp\left(-\frac{n u^2}{4} \right),
\end{equation*}
as desired.
